\documentclass[11pt,a4paper]{article}
\usepackage{../../common/preamble}

\begin{document}

\begin{center}
    {\LARGE \textbf{Количество информации. Формула Хартли. Информация и вероятность.}}\\[0.2cm]
    \rule{\textwidth}{1pt}
\end{center}

\section{Математический аппарат}

\subsection{Логико-математический синтаксис}
\begin{itemize}
    \item $\forall$ --- квантор всеобщности (для любого, для всякого)
    \item $\exists$ --- квантор существования (существует, найдется хотя бы один)
    \item $\exists!$ --- квантор существования и единственности (существует ровно один)
    \item $\land$ (\&) --- конъюнкция (логическое И)
    \item $\lor$ --- дизъюнкция (логическое ИЛИ)
    \item $\neg A$ (или $\overline{A}$) --- отрицание события/высказывания $A$
    \item $\implies$ --- импликация (логическое следование)
    \item $\iff$ --- эквивалентность (тогда и только тогда, когда)
    \item $\lfloor x \rfloor$ --- пол: наибольшее целое число, не превосходящее $x$:
    $$\lfloor x \rfloor \defeq \max \{ n \in \mathbb{Z} \mid n \le x \}$$
    \item $\lceil x \rceil$ --- потолок: наименьшее целое число, не меньшее $x$:
    $$\lceil x \rceil \defeq \min \{ n \in \mathbb{Z} \mid n \ge x \}$$
\end{itemize}

\subsection{Определение и свойства логарифма}
\begin{definitionbox}{Логарифм}
Логарифмом числа $b$ по основанию $a$ ($b > 0, a > 0, a \neq 1$) называется показатель степени $x$, в которую нужно возвести $a$, чтобы получить $b$:
$$\log_a b = x \iff a^x = b$$
\end{definitionbox}

\textbf{Свойства логарифмической функции:}
\begin{enumerate}
    \item $\log_a (x \cdot y) = \log_a x + \log_a y$
    \item $\log_a \left(\frac{x}{y}\right) = \log_a x - \log_a y$
    \item $\log_a (x^k) = k \cdot \log_a x$
    \item $\log_a \left(\frac{1}{x}\right) = -\log_a x$
    \item $\log_a b = \frac{\log_c b}{\log_c a}$
\end{enumerate}

\section{Алфавитный подход и формула Хартли}

\subsection{Интуиция и принцип аддитивности информации}
Попробуем построить математическую меру информации. Рассмотрим правильный шестигранный игральный кубик.
Если мы бросаем игральный кубик $1$ раз, мы получаем некоторую порцию информации. Если мы бросаем кубик $2$ раза ($3$ раза, $k$ раз), согласно здравому смыслу, суммарный объем полученной информации увеличится ровно в $2$ раза ($3$ раза, $k$ раз).
Это требование называют \textbf{принципом аддитивности} (информация от независимых источников должна складываться).

С другой стороны, посмотрим на пространство элементарных исходов:
\begin{itemize}
    \item При $1$ броске кубика возможно $N_1 = 6$ исходов.
    \item При $2$ бросках: $N_2 = 6 \cdot 6 = N_1 \cdot N_1 = 6^2 = 36$ исходов.
    \item При $k$ независимых бросках: $N_k = 6^k$ исходов.
\end{itemize}
Количество исходов растет \textbf{мультипликативно} (экспоненциально), а желаемая мера информации $I$ обязана расти \textbf{линейно} (аддитивно), то есть, выражаясь формальным языком, должно выполняться равенство:
$$I(N_1 \cdot N_2) = I(N_1) + I(N_2)$$

Из математического анализа известно функциональное уравнение Коши: единственная непрерывная монотонная функция, переводящая произведение аргументов в сумму значений, --- это логарифм:
$$I(N) = C \cdot \log_a N$$

\subsection{Формула Хартли (1928 г.)}
Положив нормировочную константу $C = 1$ и выбрав основание $a = 2$, получают единицу измерения \textbf{бит}:
\begin{definitionbox}{Формула Хартли}
Если дискретная система делает выбор одного из $N$ равновероятных исходов, то количество информации $I$, полученное в сообщении о реализации конкретного исхода, равно:
$$I = \log_2 N \quad [\text{бит}]$$
\end{definitionbox}

\subsection{* Другие основания и единицы информации}
Выбор основания $a = 2$ продиктован исключительно инженерным удобством двоичной схемотехники (наличие или отсутствие потенциала). С точки зрения чистой математики основание может быть любым вещественным числом $a > 1$:
$$I_a = \log_a N$$

Переход между единицами измерения осуществляется по правилу замены основания логарифма:
$$I_a = I_b \cdot \log_a b$$

\begin{table}[h!]
\centering
\begin{tabular}{@{}lll@{}}
\toprule
\textbf{Основание $a$} & \textbf{Единица измерения} & \textbf{Связь с битом} \\ \midrule
$a = 2$ & \textbf{Бит} & $\equiv 1$ бит \\
$a = e \approx 2.718$ & \textbf{Нат} & $1 \text{ нат} = \log_2 e \approx 1.443 \text{ бит}$ \\
$a = 10$ & \textbf{Децит/Хартдит} & $1 \text{ децит} = \log_2 10 \approx 3.322 \text{ бит}$ \\
$a = 3$ & \textbf{Трит} & $1 \text{ трит} = \log_2 3 \approx 1.585 \text{ бит}$ \\ \bottomrule
\end{tabular}
\end{table}

\begin{infobox}{Почему появились другие основания?}
\begin{enumerate}
    \item \textbf{Нат:} В непрерывной теории информации и статистике формулы содержат интегралы и дифференциалы. Поскольку $(\ln x)' = \frac{1}{x}$, а $(\log_2 x)' = \frac{1}{x \ln 2}$, вычисления в \textit{натах} избавляют от лишних коэффициентов $\ln 2$.
    \item \textbf{Трит:} В схемотехнике важна аппаратная сложность: число состояний системы из $m$ разрядов по $n$ уровней равно $N = n^m$, а затраты оборудования пропорциональны величине $E = n \cdot m$. Отсюда $m = \frac{\ln N}{\ln n} \implies E(n) = \ln N \cdot \frac{n}{\ln n}$. Минимум функции $f(n) = \frac{n}{\ln n}$ достигается строго при $n = e \approx 2.718$. Ближайшее натуральное число --- $3$. Троичные ЭВМ теоретически наиболее плотно используют логические элементы.
    \item \textbf{Децит:} Р.~Хартли (1928~г.) выбрал основание $10$ из-за десятичной телефонии (шаговые искатели) и таблиц Бригса. Позже А.~Тьюринг использовал эту меру для удобного сложения целых чисел вручную при взломе \guillemotleft Энигмы\guillemotright.
\end{enumerate}
\end{infobox}

\subsection{Мощность алфавита и информационная емкость}
Пусть дан алфавит $\Sigma = \{s_1, s_2, \dots, s_M\}$ мощности $|\Sigma| = M$. Сообщение представляет собой последовательность (слово) длины $L$ в данном алфавите. Общее число возможных сообщений по комбинаторным соображениям равно $N = M^L$.

По формуле Хартли:
$$I = \log_2 (M^L) = L \cdot \log_2 M = L \cdot i,$$
где $i \defeq \log_2 M$ --- \textbf{информационный вес одного символа} алфавита.

\begin{questionbox}{Дробная мера и физическая дискретизация}
Пусть алфавит состоит из трех символов ($M = 3$). Вес символа равен:
$$i = \log_2 3 \approx 1.585 \text{ бит}$$
Как трактовать нецелое число бит и как это реализуется в архитектуре компьютера?
\tcblower
Информационная емкость сообщения длины $L = 1000$ в таком алфавите равна $1000 \cdot \log_2 3 \approx 1585$ бит. 
\begin{itemize}
    \item При \textbf{посимвольном кодировании} на вычислительной технике мы обязаны выделять целое число физических триггеров (бит) на один знак: 
    $$i_{\text{физ}} = \lceil \log_2 3 \rceil = 2 \text{ бита}$$ 
    При этом избыточность составляет $2 - 1.585 = 0.415$ бит/символ (потеря $\approx 20.7\%$ памяти).
    \item Для преодоления неэффективности применяют \textbf{блочное кодирование}: кодируя группы по $k$ символов блоками бит. Например, при блоках длины 5: $3^5 = 243 \le 256 = 2^8$. Для 5 символов достаточно 8 бит, то есть $\frac{8}{5} = 1.6$ бита/символ, что сильно ближе к числу, полученному по формуле Хартли.
\end{itemize}
\end{questionbox}

\section{Вероятностный подход к измерению информации}

\subsection{Неравновероятные исходы и собственная информация события}
Вернемся к аналогии с игральной костью, но предположим, что кость со смещенным центром тяжести, и грани выпадают с неравными вероятностями $p_i$. Здравый смысл подсказывает: сообщение о выпадении крайне редкого исхода должно нести существенно больше информации, чем сообщение об исходе высоковероятном.

Какими математическими свойствами должна обладать функция меры информации $I(p)$ от вероятности события $p \in (0, 1]$? Сформулируем три базовых требования:
\begin{enumerate}
    \item \textbf{Монотонное убывание по вероятности:} чем меньше вероятность наступления события, тем выше информационная ценность сообщения о нем:
    $$\forall p_1, p_2 \in (0, 1]: \quad p_1 < p_2 \implies I(p_1) > I(p_2)$$
    \item \textbf{Нулевая информативность достоверного исхода:} наступление события с вероятностью $p = 1$ не снимает никакой неопределенности (мы знали исход заранее), следовательно, не несет новой информации:
    $$I(1) = 0$$
    \item \textbf{Аддитивность для независимых событий:} если события $A$ и $B$ статистически независимы, то вероятность их совместного свершения мультипликативна: $P(A \cap B) = P(A) \cdot P(B)$. При этом мера информации обязана складываться:
    $$I(P(A) \cdot P(B)) = I(P(A)) + I(P(B))$$
\end{enumerate}

Из математического анализа известно, что единственным классом непрерывных функций, удовлетворяющих уравнению Коши $f(x \cdot y) = f(x) + f(y)$ и условию монотонного убывания, является логарифм: $f(p) = -C \log_a p$ (где $C > 0, a > 1$). Положив $C = 1$ и $a = 2$, получаем определение:

\begin{definitionbox}{Собственная информация события}
Количество информации, содержащееся в сообщении о наступлении случайного события $E$, имеющего априорную вероятность $P(E) = p \in (0, 1]$, называется \textbf{собственной информацией события} и равно:
$$I(E) \defeq -\log_2 p = \log_2 \left(\frac{1}{p}\right) \quad [\text{бит}]$$
\end{definitionbox}

\textit{Следствие:} В частном случае ансамбля из $N$ равновероятных событий ($p = \frac{1}{N}$) формула собственной информации вырождается в формулу Хартли:
$$I(E) = -\log_2 \left(\frac{1}{N}\right) = \log_2 N$$

\subsection{Информационная энтропия Шеннона}
\begin{definitionbox}{Энтропия Шеннона}
Пусть дискретная случайная величина $X$ принимает значения из множества $\{x_1, x_2, \dots, x_n\}$ с распределением вероятностей $\{p_1, p_2, \dots, p_n\}$, причем $\sum_{i=1}^n p_i = 1$. 

\textbf{Информационной энтропией Шеннона} $H(X)$ называется математическое ожидание собственной информации исходов:
$$H(X) = \mathbb{E}[I(X)] = -\sum_{i=1}^n p_i \log_2 p_i = \sum_{i=1}^n p_i \log_2 \frac{1}{p_i}$$
\textit{(если $p_i = 0$, слагаемое принимается равным $0$, так как $\lim_{p \to 0^+} p \log_2 p = 0$)}.
\end{definitionbox}

Энтропия определяет среднюю степень априорной неопределенности системы. Количество полученной информации вычисляется через редукцию неопределенности:
$$\Delta H = H_{\text{нач}} - H_{\text{кон}}$$

\begin{questionbox}{Максимум энтропии}
При каком распределении вероятностей $\{p_i\}_{i=1}^n$ энтропия Шеннона $H(X)$ принимает максимальное значение?
\tcblower
\textbf{Ответ:} При равномерном распределении: $\forall i \in \{1, \dots, n\}: p_i = \frac{1}{n}$.\\
В этом случае неопределенность максимальна:
$$H_{\max} = -\sum_{i=1}^n \frac{1}{n} \log_2 \frac{1}{n} = \log_2 n$$
\textit{Упражнение:} докажите аналитически для $n=2$ через поиск критических точек функции $f(p) = -p \log_2 p - (1-p) \log_2 (1-p)$ с помощью производной.
\end{questionbox}

\section{Практика}

\subsection{ЕГЭ №11: Равномерное кодирование текстовой информации}

\begin{taskbox}{Пароли с дополнительными сведениями}
При регистрации в компьютерной системе каждому пользователю выдаётся пароль, состоящий из 15 символов и содержащий только символы из 12-символьного набора: \texttt{А, В, C, D, Е, F, G, H, К, L, M, N}. 

В базе данных для хранения сведений о каждом пользователе отведено одинаковое и минимально возможное целое число байт. При этом используют посимвольное кодирование паролей, все символы кодируют одинаковым и минимально возможным количеством бит. 

Кроме собственно пароля, для каждого пользователя в системе хранятся дополнительные сведения, для чего отведено 12 байт на одного пользователя.

Определите объём памяти (в байтах), необходимый для хранения сведений о 50 пользователях. В ответе запишите только целое число.
\end{taskbox}

\textbf{Решение:}
\begin{enumerate}
    \item Мощность алфавита пароля: $M = 12$. 
    По изученному выше вес одного символа $i$ обязан удовлетворять равенству:
    $$i = \lceil \log_2 12 \rceil = 4 \text{ бит}$$
    \item Информационный объем одного пароля длины $L = 15$ символов:
    $$I_{\text{пароль}} = L \cdot i = 15 \cdot 4 = 60 \text{ бит}$$
    \item На уровне базы данных под пароль резервируется целое число байт:
    $$V_{\text{пароль}} = \left\lceil \frac{60}{8} \right\rceil = 8 \text{ байт}$$
    \item Суммарный объем данных на одного пользователя складывается из объема пароля и фиксированного объема дополнительных полей:
    $$V_{\text{пользователь}} = V_{\text{пароль}} + V_{\text{доп}} = 8 + 12 = 20 \text{ байт}$$
    \item Итоговый объем памяти для хранения данных о 50 пользователях:
    $$V_{\text{общ}} = 50 \cdot V_{\text{пользователь}} = 50 \cdot 20 = 1000 \text{ байт}$$
\end{enumerate}
\textbf{Ответ:} 1000

\subsection{ЕГЭ №7: Равномерное кодирование графической и звуковой информации}

\begin{taskbox}{Кодирование изображений}
Прибор автоматической фиксации нарушений делает цветные фотографии размером $1024 \times 768$ пикселей, используя палитру из 4096 цветов. Для пакетной передачи снимки группируются по 256 штук. Определите размер одного пакета фотографий в Мбайт.
\end{taskbox}

\textbf{Решение:}
\begin{enumerate}
    \item Глубина цвета (информационный вес одного пикселя):
    $$M = 4096 \implies i = \log_2 4096 = 12 \text{ бит}$$
    \item Информационный объем одного кадра:
    $$I_{\text{кадр}} = 1024 \cdot 768 \cdot 12 \text{ бит} = 2^{10} \cdot (3 \cdot 2^8) \cdot 12 \text{ бит}$$
    \item Переведем объем кадра в килобайты:
    $$I_{\text{кадр}} = \frac{2^{10} \cdot 768 \cdot 12}{8 \cdot 1024} = \frac{768 \cdot 12}{8} = 768 \cdot 1.5 = 1152 \text{ Кбайт}$$
    \item Объем пакета из 256 фотографий:
    $$V_{\text{пакет}} = 256 \cdot 1152 \text{ Кбайт} = \frac{2^8 \cdot 1152}{2^{10}} \text{ Мбайт} = 288 \text{ Мбайт}$$
\end{enumerate}
\textbf{Ответ:} 288

\begin{taskbox}{Кодирование звуков}
Производится двухканальная (стерео) звукозапись с частотой дискретизации 64 кГц и 24-битным разрешением. В результате был получен файл размером 220 Мбайт без учета заголовка и сжатия данных. Определите длительность звукозаписи в минутах. В качестве ответа укажите ближайшее к полученному времени целое число.
\end{taskbox}

\textbf{Решение:}
\begin{enumerate}
    \item Параметры оцифровки:
    \begin{itemize}
        \item Число каналов: $k = 2$ (стерео).
        \item Частота дискретизации: $f = 64\text{ кГц} = 64\,000\text{ Гц}$.
        \item Разрешение: $i = 24\text{ бита}$.
    \end{itemize}
    \item Битрейт аудиопотока:
    $$B = k \cdot f \cdot i = 2 \cdot 64\,000 \cdot 24 = 3\,072\,000 \text{ бит/с}$$
    \item Длительность звукозаписи $T$:
    $$T = \frac{I}{B} = \frac{220 \cdot 2^{23}}{3\,072\,000} = \approx 600.1 \text{ с} \approx 10 \text{ мин}$$
\end{enumerate}
\textbf{Ответ:} 10

\end{document}